\chapter{DEFINICIONES Y ABREVIATURAS}
\label{cap:definiciones}

En este capítulo se recogen los términos técnicos y los acrónimos que aparecen
a lo largo de la memoria, ordenados alfabéticamente.

\section{Definiciones}

\begin{description}          
  
  \item[\gldestino{alucinacion}Alucinación] Respuesta de un modelo de lenguaje que resulta
        lingüísticamente plausible pero factualmente incorrecta. Es el riesgo que
        motiva el uso de una arquitectura RAG en este trabajo: un modelo sin
        fuentes puede inventar asignaturas o planes de estudio que no existen.
        
  \item[\gldestino{autoatención}Autoatención (\textit{self-attention})] Mecanismo central de la
        arquitectura Transformer que calcula la representación de cada palabra
        ponderando la información del resto de palabras de la secuencia según su
        relevancia, sin importar la distancia que las separa.
        
  \item[\gldestino{base de datos vectorial}Base de datos vectorial] Sistema de almacenamiento especializado en
        guardar incrustaciones y recuperar las más parecidas a una consulta según
        su similitud, en lugar de por coincidencia exacta de palabras.
        
  \item[\gldestino{colección documental}Colección documental] Conjunto de textos, extraídos y depurados, que
        constituye la única fuente de conocimiento del sistema. En este trabajo la
        forman las guías docentes, los planes de estudio y las salidas
        profesionales de la EPSJ.
        
  \item[\gldestino{corpus}Corpus] Colección de textos que se utiliza para entrenar o evaluar
        un modelo de lenguaje. En este trabajo, el término se usa también para la
        colección documental extraída de la web de la EPSJ.
        
  \item[\gldestino{exhaustividad}Exhaustividad (\textit{Recall@K})] \textcolor{borrador2}{Métrica de
        recuperación que mide qué proporción de los fragmentos que hacían falta
        para responder aparecen entre los $K$ primeros resultados. No juzga la
        respuesta, sino el material con el que el modelo va a contar para
        redactarla. Su definición formal es la Ecuación~\ref{eq:recall_at_k}.}

  \item[\gldestino{fidelidad del índice}Fidelidad del índice] \textcolor{borrador2}{Grado en que una base de
        datos vectorial devuelve exactamente los mismos vecinos que la búsqueda
        exacta sobre los mismos vectores. Es una propiedad de la herramienta de
        almacenamiento, y se comprueba comparándola con una búsqueda por fuerza
        bruta que sirve de referencia.}

  \item[\gldestino{fidelidad de la respuesta}Fidelidad de la respuesta (\textit{faithfulness})]
        \textcolor{borrador2}{Grado en que lo que afirma una respuesta generada
        está respaldado por los fragmentos que se le entregaron como contexto.
        Es una propiedad del modelo de lenguaje, no del almacenamiento, y es lo
        que se mide para detectar alucinaciones. Conviene no confundirla con la
        fidelidad del índice: son medidas distintas, sobre piezas distintas del
        sistema y en momentos distintos.}

  \item[\gldestino{chunking}Fragmentación (\textit{chunking})] Proceso de dividir cada documento
        de la colección en fragmentos de un tamaño manejable para el modelo de
        incrustaciones, procurando que cada uno conserve sentido por sí solo.
        
  \item[\gldestino{fragmento}Fragmento (\textit{chunk})] Cada una de las piezas de texto
        resultantes de la fragmentación. Es la unidad mínima que el sistema
        recupera y entrega al modelo de lenguaje como contexto.

  \item[\gldestino{generación aumentada por recuperación}Generación aumentada por recuperación] Arquitectura que, antes
        de responder, busca en una colección documental los fragmentos relevantes
        para la pregunta y se los entrega al modelo de lenguaje como contexto, de
        modo que su respuesta se apoye en documentos reales y verificables. Véase
        también RAG.

  \item[\gldestino{guía docente}Guía docente] Documento oficial que publica la universidad para cada
        asignatura, con su resumen, sus requisitos previos y su temario.
        
  \item[\gldestino{incrustacion}Incrustación (\textit{embedding})] Representación de un texto como un
        vector numérico, construida de forma que los textos con significados
        parecidos queden próximos entre sí. Es lo que permite buscar por
        significado y no solo por palabras literales.
        
  \item[\gldestino{Minería de textos y datos}Minería de textos y datos] \textcolor{borrador2}{Toda técnica analítica
        automatizada destinada a analizar textos y datos en formato digital con
        el fin de generar información, ya sean pautas, tendencias o
        correlaciones. La definición es la que recoge el artículo 66.1 del Real
        Decreto-ley 24/2021\cite{rdl24_2021}, y se adopta por ser la que fija su
        alcance legal: es el supuesto bajo el que se ampara la extracción
        automatizada realizada en este trabajo (Sección~\ref{secc:marcolegal}).}

  \item[\gldestino{integración continua}Integración Continua] Práctica de ingeniería que ejecuta de forma
        automática la batería de pruebas cada vez que se sube código al
        repositorio, de modo que un fallo se detecta en el momento de introducirlo.
        
  \item[\gldestino{Kanban}Kanban] Método ágil de gestión del trabajo basado en la
        visualización del flujo de tareas en un tablero y en limitar el trabajo
        en curso. Es la metodología con la que se ha gestionado este proyecto.
        
  \item[\gldestino{modelo de lenguaje grande}Modelo de lenguaje grande (LLM)] Sistema computacional entrenado
        sobre grandes volúmenes de texto para predecir la palabra más probable a
        partir del contexto anterior, capacidad de la que emergen la generación y
        la comprensión de lenguaje natural que exhiben estos modelos.
        
  \item[\gldestino{rango recíproco medio}Rango recíproco medio (\textit{Mean Reciprocal Rank}, MRR)]
        \textcolor{borrador2}{Métrica de recuperación que premia colocar el
        primer resultado correcto lo más arriba posible, y no solo colocarlo:
        toma el inverso de la posición en que aparece y lo promedia sobre todas
        las preguntas. Su definición formal es la Ecuación~\ref{eq:mrr}.}

  \item[\gldestino{spider}Rastreador web (\textit{spider})] Programa que recorre de forma
        automática las páginas de un sitio web siguiendo sus enlaces, para extraer
        de ellas datos estructurados.

  \item[\gldestino{señal semántica}Señal semántica] Parte del significado de un
        texto que su incrustación consigue representar. Un fragmento centrado en un
        único asunto produce un vector que apunta con claridad hacia él, uno que
        mezcla asuntos distintos produce un vector intermedio entre todos ellos, que
        no representa bien a ninguno. De ahí la expresión <<diluir la señal
        semántica>>.
        
  \item[\gldestino{slug}\textit{Slug}] Parte final de una dirección web que identifica el
        recurso concreto.
        
  \item[\gldestino{test de Holland}Test de Holland] Instrumento de
        orientación vocacional que clasifica los intereses profesionales de una
        persona en seis tipos (realista, investigador, artístico, social,
        emprendedor, convencional) y los compara con perfiles de puestos o
        estudios afines, para sugerir cuáles encajan mejor con la persona
        evaluada.
        
  \item[\gldestino{token}\textit{Token}] Unidad mínima de texto con la que opera un modelo de
        lenguaje (una palabra, una parte de palabra o un signo de puntuación). Los
        modelos no procesan caracteres ni palabras completas, sino secuencias de
        \textit{tokens}.
        
  \item[\gldestino{trabajo en curso}Trabajo en curso (\textit{Work In Progress}, WIP)] Número de tareas
        abiertas simultáneamente. Limitarlo es la regla central de Kanban: obliga a
        terminar antes de empezar algo nuevo.
        
  \item[\gldestino{transformador}Transformador (\textit{Transformer})] Arquitectura de red neuronal
        basada en el mecanismo de autoatención, que prescinde de la recurrencia y
        permite procesar las secuencias en paralelo. Es la base arquitectónica de
        los modelos de lenguaje grandes actuales.
        
  \item[\gldestino{ventana de contexto}Ventana de contexto] Cantidad máxima de texto,
        medida en \textit{tokens}, que un modelo puede procesar en una sola
        operación. El texto que excede ese límite se descarta sin aviso, por lo que
        la ventana del modelo de incrustaciones es la que fija el tamaño máximo
        admisible de un fragmento.
        
\end{description}

\section{Abreviaturas y acrónimos}

\begin{description}          
  \item[\gldestino{ADR}ADR] Registro de decisión de arquitectura (\textit{Architecture
          Decision Record}). Documento breve que deja constancia de una decisión
        técnica, sus alternativas y sus motivos.
  \item[\gldestino{CI}CI] Integración continua (\textit{Continuous Integration}).
  \item[\gldestino{DQA}DQA] Evaluación de calidad de datos (\textit{Data Quality
          Assessment}). Registro con el mismo rigor que un ADR, pero para
        anomalías de calidad de datos en lugar de decisiones de arquitectura:
        qué se observó en la fuente, dónde y qué tratamiento recibió.
  \item[\gldestino{ECTS}ECTS] Sistema europeo de transferencia y acumulación de créditos.
  \item[\gldestino{EPSJ}EPSJ] Escuela Politécnica Superior de Jaén.
  \item[\gldestino{HTML}HTML] Lenguaje de marcado de hipertexto, en el que se publican las
        páginas web de las que se extraen los datos.
  \item[\gldestino{IA}IA] Inteligencia Artificial.
  \item[\gldestino{LLM}LLM] Modelo de lenguaje grande (\textit{Large Language Model}).
  \item[\gldestino{MoSCoW}MoSCoW] Técnica de priorización que clasifica las tareas en
        \textit{Must}, \textit{Should}, \textit{Could} y \textit{Won't}.
  \item[\gldestino{MRR}MRR] \textcolor{borrador2}{Rango recíproco medio (\textit{Mean
          Reciprocal Rank}).}
  \item[\gldestino{NLG}NLG] Generación de lenguaje natural (\textit{Natural Language
          Generation}), la vertiente del PLN dedicada a producir texto.
  \item[\gldestino{NLP}NLP] Procesamiento del lenguaje natural (\textit{Natural Language
          Processing}); véase PLN.
  \item[\gldestino{PLN}PLN] Procesamiento del lenguaje natural.
  \item[\gldestino{RAG}RAG] Generación aumentada por recuperación (\textit{Retrieval-Augmented Generation}).
  \item[\gldestino{TFG}TFG] Trabajo Fin de Grado.
  \item[\gldestino{TGBA}TGBA] Tasa Global Bruta de Abandono.
  \item[\gldestino{TGNA}TGNA] Tasa Global Neta de Abandono.
  \item[\gldestino{UJA}UJA] Universidad de Jaén.
  \item[\gldestino{URL}URL] Localizador uniforme de recursos; la dirección de una página web.
  \item[\gldestino{WIP}WIP] Trabajo en curso (\textit{Work In Progress}).
\end{description}